% Sample LaTeX file for PMLR single-column format
% This is the standard format for most PMLR proceedings

\documentclass[pmlr]{jmlr} % This is the PMLR single-column format

% Recommended packages
\usepackage{booktabs} % For better tables
\usepackage{graphicx} % For including figures
\usepackage{url}      % For URLs
\usepackage{amsmath}  % For mathematical notation
\usepackage{algorithm2e} % For algorithms (optional)

% Proceedings information
\jmlrvolume{123}
\jmlryear{2024}
\jmlrworkshop{Workshop on Machine Learning Theory}

% Title and author information
\title{Sample Paper Title: A Comprehensive Approach to Machine Learning}

% Use \Name{Author Name} to specify the name.
% If the surname contains spaces, enclose the surname
% in braces, e.g. \Name{John {Smith Jones}} similarly
% if the name has a "von" part, e.g \Name{Jane {de Winter}}.
% If the first letter in the forenames is a diacritic
% enclose the diacritic in braces, e.g. \Name{{\'E}louise Smith}

% Authors with diverse names demonstrating proper formatting:
\author{\Name{Mar\'ia Garc\'ia-L\'opez} \Email{maria.garcia@example.edu}\\
  \addr Department of Computer Science\\
  Technical University of Madrid\\
  Madrid, Spain
  \AND
  \Name{Li Wei} \Email{li.wei@example.edu.cn}\\
  \addr Institute of Computing Technology\\
  Chinese Academy of Sciences\\
  Beijing, China
  \AND
  \Name{Jamal O'Brien-Williams} \Email{jamal.obrien@example.ie}\\
  \addr School of Computer Science\\
  Trinity College Dublin\\
  Dublin, Ireland
  \AND
  \Name{Nguy\~{e}n Th\d{i} Minh} \Email{nguyen.minh@example.edu.vn}\\
  \addr Faculty of Information Technology\\
  University of Science, VNU-HCM\\
  Ho Chi Minh City, Vietnam
}

\editor{Sarah Johnson and Robert Chen}

\begin{document}

\maketitle

\begin{abstract}
This is a sample abstract for a PMLR paper in single-column format. 
The abstract should be a brief summary of your paper, typically 
100-200 words. It should clearly state the problem, your approach, 
and the main results. Mathematical notation can be included, such as 
$f(x) = \sum_{i=1}^{n} w_i x_i$, and the abstract can span multiple 
lines as needed.
\end{abstract}

\begin{keywords}
Machine Learning, Sample Paper, PMLR Format, Single Column
\end{keywords}

\section{Introduction}

This is a sample PMLR paper in single-column format. The single-column
format is the standard format for many PMLR proceedings. This format
uses the \texttt{jmlr} class with the \texttt{pmlr} option, which
automatically sets the header to read ``Proceedings of Machine Learning
Research''.

You can include citations using standard \LaTeX{} citation commands,
such as \citet{example2023} or \citep{example2023}. Mathematical equations can be displayed:

\begin{equation}
  \label{eq:example}
  f(x) = \int_{-\infty}^{\infty} e^{-x^2} dx = \sqrt{\pi}
\end{equation}

\section{Methods}

Here you would describe your methods. You can include figures using
the standard \texttt{figure} environment:

\begin{figure}[htbp]
  \centering
  % Uncomment the line below when you have an actual figure
  % \includegraphics[width=0.8\columnwidth]{figure.pdf}
  \caption{This is a sample figure caption. Figures should be
    referenced in the text using \texttt{Figure~\ref{fig:sample}}.}
  \label{fig:sample}
\end{figure}

\subsection{Algorithms}

You can also include algorithms. Here's a simple example:

\begin{algorithm}[htbp]
  \SetAlgoLined
  \KwIn{Training data $(x_1, y_1), \ldots, (x_n, y_n)$}
  \KwOut{Model parameters $\theta$}
  Initialize $\theta$ randomly\;
  \While{not converged}{
    Compute gradient $\nabla L(\theta)$\;
    Update $\theta \leftarrow \theta - \eta \nabla L(\theta)$\;
  }
  \Return $\theta$\;
  \caption{Sample gradient descent algorithm}
  \label{alg:sample}
\end{algorithm}

\section{Results}

Present your results here. Tables can be formatted using the
\texttt{booktabs} package:

\begin{table}[htbp]
  \centering
  \caption{Sample results table}
  \label{tab:results}
  \begin{tabular}{lcc}
    \toprule
    Method & Accuracy & Training Time (s) \\
    \midrule
    Baseline & 85.2\% & 120 \\
    Our Method & 92.1\% & 95 \\
    \bottomrule
  \end{tabular}
\end{table}

\section{Conclusion}

Summarize your findings and their implications here.

\section*{Acknowledgments}

Acknowledge funding sources and collaborators here. Note that this
section is unnumbered (using \texttt{section*}).

% Bibliography
% PMLR uses author-year citations by default
% Use \citet{} for textual citations: "Smith and Jones (2023) showed..."
% Use \citep{} for parenthetical citations: "as shown previously (Smith and Jones, 2023)"

% You can use BibTeX (recommended) or write entries manually as shown below
% For BibTeX, use: \bibliography{yourbibfile}

\begin{thebibliography}{10}

\bibitem[Author and Collaborator(2023)]{example2023}
Author, A. and Collaborator, B.
\newblock Sample Paper Title.
\newblock \emph{Proceedings of Machine Learning Research}, 123:456--789, 2023.

\end{thebibliography}

% Appendix (optional)
\appendix
\section{Additional Details}

Any supplementary material, proofs, or additional details can be
included in an appendix.

\end{document}

